\section{Method}

\subsection{Benchmark objective}

The benchmark compares forecasting model families for weekly influenza
hospitalization-rate prediction under a common data-processing, fitting, and
evaluation protocol. The objective is twofold:
\begin{enumerate}[leftmargin=1.5em]
\item quantify point-forecast performance across forecasting baselines,
hand-specified epidemic models, and discovered epidemic structures; and
\item determine whether constrained structure discovery provides useful
age-specific latent or observation structures beyond strong baselines.
\end{enumerate}

All model selection and hyperparameter selection use only training,
validation, or rolling-validation evidence. Held-out test metrics are computed
only after the model family, structure, or baseline configuration has been
selected.

\subsection{Data processing and series construction}

The raw FluSurv-NET CSV is read with \texttt{skiprows=2} to discard metadata
lines above the table. Rows with missing \texttt{WEEK} or \texttt{YEAR.1} are
removed, the remaining rows are sorted chronologically by year and week, and a
continuous integer time index is created. The primary series is the
\emph{Entire Network} / \emph{Overall} weekly hospitalization rate. The frozen
run also evaluates five age-stratified series:
\texttt{0--4 yr}, \texttt{5--17 yr}, \texttt{18--49 yr},
\texttt{50--64 yr}, and \texttt{$\geq$65 yr}.

\subsection{Forecasting baselines}

The discovery-ablation configuration adds opt-in non-epidemic forecasting
baselines. Persistence and rolling-mean baselines test whether recent observed
rates are already sufficient. \texttt{arima\_auto\_small} chooses among a
small fixed set of ARIMA orders by validation MAE and then refits on
train-plus-validation for test forecasting. \texttt{lagged\_ridge} and
\texttt{lagged\_gradient\_boosting} use lagged observations, linear time, and
seasonal sine/cosine features with recursive multi-step forecasts. These
baselines are opt-in; the historical default benchmark remains the original
six epidemic model families.

\subsection{Manual epidemic baselines}

The hand-specified epidemic baselines include deterministic SEIR,
probabilistic SEIR, hospitalized SEIHR, delayed-observation SEIR, and
fractional SEIR variants. All epidemic states are normalized, and fitting uses
explicit penalties for negative states and approximate mass-conservation
violations. The delayed-observation SEIR baseline selects its delay using
validation evidence only.

\subsection{Constrained structure discovery}

The discovery module searches over a bounded grammar rather than arbitrary
symbolic dynamics. A structure specification contains:
\begin{itemize}[leftmargin=1.5em]
\item a compartment template from
\{\texttt{SIR}, \texttt{SEIR}, \texttt{SEIRS}, \texttt{SEIHR}, \texttt{SEIAR}\},
\item a fractional on/off flag,
\item an observation map from
\{\texttt{I}, \texttt{H}, \texttt{I+H}, \texttt{delayed\_I}\}, and
\item a delay parameter in $\{0,1,2,3\}$ when the observation map is
\texttt{delayed\_I}.
\end{itemize}

Validity rules enforce that infection originates from $S$, infectious-like
states eventually reach recovery, no isolated compartments exist, and
observation semantics are compatible with the latent template. For example,
\texttt{obs=H} and \texttt{obs=I+H} require an $H$ compartment, while
\texttt{obs=delayed\_I} requires an $I$ compartment and a valid delay.

\subsection{Discovery score policies and ablations}

The constrained discovery model uses a stability-aware score that combines
multi-split validation error, rolling-origin stability, complexity penalties,
and the configured age prior. The frozen ablation package compares this search
against five same-grammar alternatives:
\begin{itemize}[leftmargin=1.5em]
\item \texttt{random\_structure\_discovery}, which evaluates a seeded random
candidate budget;
\item \texttt{exhaustive\_structure\_discovery}, which evaluates the valid
grammar universe subject to an explicit candidate guard;
\item \texttt{validation\_only\_structure\_selection}, whose score is exactly
validation MAE;
\item \texttt{no\_observation\_search\_discovery}, which restricts the
candidate universe to \texttt{obs=I}; and
\item \texttt{no\_stability\_discovery}, which removes the rolling stability
penalty from the score while retaining other terms.
\end{itemize}

All discovery leaderboards include score-policy audit columns so that the
selected structure can be traced back to the evidence and penalties used.

\subsection{Post-hoc paired rolling comparison}

After model selection is fixed, the analysis aligns rolling-origin forecasts
by series, horizon, and target week. For each challenger and the constrained
discovery reference, it reports paired mean absolute-error differences,
bootstrap confidence intervals, and reference win rates. These paired
comparisons summarize post-hoc evidence and are not used for model selection.
